\documentclass[11pt]{amsart}
\usepackage{amsmath,amssymb,amsthm}
\usepackage[margin=1.1in]{geometry}
\usepackage{booktabs,array}
\usepackage[colorlinks=true,linkcolor=blue,citecolor=blue,urlcolor=blue]{hyperref}
\hypersetup{pdftitle={The Hodge conjecture for Fermat fourfolds of odd degree at most 199}, pdfauthor={Rifat Jumagulov}}
\newtheorem{theorem}{Theorem}[section]
\newtheorem{lemma}[theorem]{Lemma}
\newtheorem{proposition}[theorem]{Proposition}
\newtheorem{corollary}[theorem]{Corollary}
\theoremstyle{remark}
\newtheorem{remark}[theorem]{Remark}
\newcommand{\Q}{\mathbb{Q}}
\newcommand{\Z}{\mathbb{Z}}
\newcommand{\F}{\mathbb{F}}
\newcommand{\PP}{\mathbb{P}}
\newcommand{\Ds}{\mathfrak{D}}
\newcommand{\fr}[1]{\langle #1\rangle}
\newcommand{\V}{V}
\DeclareMathOperator{\Gal}{Gal}
\newcommand{\zt}{\zeta}
\newcommand{\C}{\mathbb{C}}
\newcommand{\Qbar}{\overline{\Q}}

\title[Hodge classes on odd-degree Fermat fourfolds]{The Hodge conjecture for Fermat fourfolds
of odd degree at most 199}
\author{Rifat Jumagulov}
\email{jum.rifm@gmail.com}
\subjclass[2020]{Primary 14C30, 14C25; Secondary 14J70, 11T24}
\keywords{Hodge conjecture, Fermat varieties, algebraic cycles, Jacobi sums, computer-assisted
proof}
\date{July 2026}

\begin{document}
\raggedbottom\vfuzz=3pt
\begin{abstract}
Let $X^4_m=\{x_0^m+\dots+x_5^m=0\}\subset\PP^5$ be the Fermat fourfold of degree $m$. We prove
the Hodge conjecture for $X^4_m$ for every odd $m\le199$.

The proof introduces three closure mechanisms for Hodge $(2,2)$-characters. First, if the
character multiset splits into two zero-sum triples, then the corresponding rational Hodge
block is transported from a $(1,1)$-substructure of a product of Fermat curves and is therefore
algebraic. Second, algebraicity follows for characters that, after adjoining two vanishing
pairs, decompose into an Aoki standard sextuple and a grade-$2$ Hodge quadruple. Third, the
exceptional class at $m=33$ has a lift to level $66$ that is quasi-decomposable; its
algebraicity then descends along the finite morphism $X^4_{66}\to X^4_{33}$.

Together with the known decomposable, quasi-decomposable, and standard cases, these criteria
cover all odd degrees $m\le199$. This is established by a complete computer-assisted census,
accompanied by a completeness proof for the enumeration, re-verifiable per-orbit certificates
(positive witnesses and, for the beyond-machinery orbits, the negative screenings), and an
algorithmically independent reimplementation reproducing the census in full.
Finally, an exact Jacobi-sum computation at $p=67$ shows that no cycle defined over
$\Q(\zt_{33})$ projects nontrivially onto the exceptional $m=33$ block. More generally, for every finite
extension of $\Q(\zt_{33})$ over which a certifying cycle is defined, every residue degree
above $67$ is divisible by $6$.
\end{abstract}
\maketitle

\section{Introduction and main results}\label{sec:intro}
Let $X^4_m\subset\PP^5$ be the Fermat fourfold $x_0^m+\dots+x_5^m=0$, $G=\mu_m^6/\mu_m$, and for a
character $a=(a_0,\dots,a_5)\in(\Z/m)^6$ with $\sum a_i\equiv0$, $a_i\neq0$, let $\V(a)\subset
H^4_{\mathrm{prim}}$ be the ($1$-dimensional) eigenspace; $a$ is a Hodge $(2,2)$ character iff
$\sum_i\fr{ta_i}=3m$ for all $t\in(\Z/m)^\times$, where $\fr{x}\in\{1,\dots,m-1\}$ is the residue
(Shioda \cite{Shioda79}). The known algebraicity machinery is: (i)~the \emph{induced structure}
(Shioda \cite{Shioda79}, Shioda--Katsura \cite{SK}), which transports algebraicity through
decompositions---da Silva's conditions $(P1)/(P2)$ are quasi-decomposability in the semigroup $M_m$
\cite{daSilva}; and (ii)~\emph{standard cycles} (Aoki \cite{Aoki87}). da Silva \cite[Prop.~3.4 (= Prop.~3.6 of the arXiv version)]{daSilva}
found the first Hodge class beyond both at $m=33$ and proposed a candidate cycle $[W]$ (Question~1).
(His closures --- $m\le20$ [Thm.~2.7], $m$ prime or $m=4$ [Thm.~2.6], $m=p^2$ [Thm.~2.8],
$m=21,27$ [Thm.~3.3], and fourfolds of degree $m\le100$ coprime to $6$ [Prop.~3.1], numbering of
the published version \cite{daSilva} --- leave the composite non-square degrees with $3\mid m$, beyond
his $21,27$, untouched; every class closed below lies there.) In summary:
\begin{center}\footnotesize
\begin{tabular}{ll>{\raggedright\arraybackslash}p{4.9cm}}
\toprule
mechanism & source & covers\\
\midrule
induced structure & Shioda \cite{Shioda79}; Shioda--Katsura \cite{SK} & decomposable, quasi-decomposable\\
standard cycles; lattice calculus & Aoki \cite{Aoki87,Aoki00} & standard characters; $\nu\in S_m$; Thm.~0.1 degree forms\\
computational closures & da Silva \cite{daSilva} & $m\le20$; $m$ prime or $4$; $p^2$; $21,27$; $m\le100$ coprime to $6$\\
Theorems~\ref{thm:A}--\ref{thm:App} $+$ census & this paper & the residual odd orbits through $199$\\
\bottomrule
\end{tabular}
\end{center}
We work in his semigroup formalism and prove:

\medskip
\noindent\emph{Notation.} Throughout $m\ge3$ ($m\le2$ gives a hyperplane or a smooth quadric,
where the conjecture is classical). A multiset $a$ of nonzero residues with $\sum a_i\equiv0\pmod m$ has
\emph{grade} $g$ if $\sum_i\fr{ta_i}=gm$ for \emph{every} $t\in(\Z/m)^\times$ (the constancy is the
Hodge condition: grade-$3$ sextuples are the Hodge $(2,2)$ characters of $X^4_m$, grade-$2$ quadruples
the $(1,1)$ characters of the surface $X^2_m$); $\sim$ is equality of multisets up to coordinate
permutation, and the \emph{content} of a class $a=(a_0,\dots,a_5)$ is $\gcd(a_0,\dots,a_5,m)$
(content-$e$ classes are
inflations from level $m/e$). $\mathrm{claim}_m(a)$ (or $\mathrm{claim}(a)$ when $m$ is clear) is Aoki's
\textsc{claim}$(a)$: $\V(a)$ is spanned by algebraic cycle classes \cite[Introduction]{Aoki87} ---
equivalently, some \emph{rational} algebraic cycle class has nonzero $\V(a)$-component,
$\omega_a(Z)\neq0$. Three levels are kept distinct throughout: eigenlines are indexed by
\emph{ordered} tuples ($\V(b)\neq\V(b')$ for distinct ordered characters even when they agree as
multisets; coordinate permutations are realized by automorphisms of $X^4_m$, which is why claim is
invariant under them, and it is likewise a Galois-orbit-level property); $M(a):=\bigoplus_b\V(b)$
is the rational Galois block --- the sum over the distinct \emph{ordered} conjugates $b=ta$,
$t\in(\Z/m)^\times$, their number being the index of the ordered stabilizer
$\{t:ta=a\text{ coordinatewise}\}$, \emph{not} the number of multiset classes; and the census
classifies sorted representatives modulo both actions --- legitimate since claim is invariant
under both, but never identified with an eigenline count. $\omega_a(Z)$ denotes the
$\V(a)$-component of the cycle class of $Z$
(Aoki's averaged projector \cite[Introduction]{Aoki87}). $*$ is Aoki's juxtaposition --- concatenation of the
entry tuples; characters of $X^r_m$ and $X^s_m$ concatenate to one of $X^{r+s+2}_m$ ($n=r+s+2$ in
Aoki's indexing by cohomological degree), grades adding \cite[\S1]{Aoki87} --- and, for even $s$, $\Ds^s_m$ his set of
$(s+2)$-entry characters that are coordinate permutations of $(a_0,-a_0,\dots,a_{s/2},-a_{s/2})$
($s/2+1$ vanishing pairs), represented by linear cycles \cite[\S1, Thm.~1-1]{Aoki87}; only $\Ds^2_m$
is used below. $M_m$ is da Silva's additive semigroup of solutions $(x_1,\dots,x_{m-1};y)$, $y>0$, of
$\sum_i\fr{ti}\,x_i=my$ for all $t$ (multiplicities intrinsic), $M_m(g)$ its grade-$g$ part; a class is
\emph{decomposable} if it splits in $M_m$ and \emph{quasi-decomposable} if it splits after adjoining
one element of $M_m(1)$ --- da Silva's $(P1)/(P2)$ \cite[Def.~2.4]{daSilva}. $\#$ is his pairing: for
$\beta=(b_0,\dots,b_{r+1})$, $\gamma=(c_0,\dots,c_{s+1})$ with nonzero entries and
$b_{r+1}+c_{s+1}\equiv0$ (the set $U^{r,s}_m$), $\beta\#\gamma=(b_0,\dots,b_r,c_0,\dots,c_s)$
\cite[before Cor.~2.3]{daSilva}.
$\sigma_{p,x}$ denotes the standard characters (Proposition~\ref{prop:D}). $u(a,p)$ is the scalar by
which geometric Frobenius at a split prime $p$ acts on $\V(a)(n/2)$ --- the Tate-normalized Jacobi sum
$j(a)/p^{n/2}$ (the Jacobi-sum count of diagonal hypersurfaces, \cite[pp.~500--502]{Weil49}; \cite{Weil52}) --- so a divisorial class has $u=1$ at split primes whose
Frobenius fixes the class of a certifying divisor (the second $m=39$ class of
\S\ref{sec:A} calibrates this in-text). Three symbols are kept
apart: $u(a,p)$ above is the \emph{local} normalized scalar; $\chi_a$ is the associated
\emph{global} Jacobi-sum Hecke character \emph{in its Tate-normalized form}: Weil's
unnormalized character $J_a$ \cite{Weil52} has $|J_a(\mathfrak p)|=N\mathfrak p^{\,n/2}$ in
every embedding ($n$ the cohomological degree of the ambient Fermat variety), so $J_a$ itself is
\emph{not} of finite order, and $\chi_a:=J_a/N^{n/2}$ is its finite-order part --- every use of
$\chi_a$ in this paper is of the normalized character, whose value at a split prime is the
local scalar $u(a,p)$; it is multiplicative under juxtaposition,
$\chi_{\alpha*\beta}=\chi_\alpha\chi_\beta$ (the normalizing powers add along the
weights; the machine-verified instances are the displayed $u$-identities); and
$\nu(a)\in\Z^{m-1}$ is Aoki's
\emph{multiplicity vector}, additive under juxtaposition,
$\nu(\alpha*\beta)=\nu(\alpha)+\nu(\beta)$ (Aoki writes $u$ for $\nu$). $S_m$ denotes
Aoki's lattice: the subgroup of $\Z^{m-1}$
generated by the $\nu(\delta)$ over the vanishing pairs $\delta=\{k,m-k\}$ and the
$\nu(\sigma_{p,x})$ over Aoki's standard elements for all primes $p\mid m$ --- for odd $p$ the AP
family of Proposition~\ref{prop:D}, at even $m$ also Aoki's $p=2$ family
$\sigma_{2,x}=(x,\,x+\frac m2,\,m-2x,\,\frac m2)$ \cite[\S2]{Aoki00} --- the ``standard
calculus''; this generating family is verbatim the generator set of Aoki's lattice $S_m$
(Proposition~\ref{prop:ident}), with membership decided exactly by integer linear
algebra (Hermite normal form), and $B^{4}_m$ the grade-$3$ Hodge-character set
of the Fermat fourfold at level $m$; in lattice statements we freely identify $a$ with $\nu(a)$
($a\in B^4_m\smallsetminus S_m$ means $\nu(a)\notin S_m$ --- a statement about the formal
calculus, not about the character group). The witness $w=(1,4,16,22,25,31)$ below is, up to coordinate permutation, the Galois
conjugate $7w$ of da Silva's displayed
representative $(7,10,13,19,22,28)\sim7w$ of \cite[Prop.~3.4 (= Prop.~3.6 of the arXiv version)]{daSilva}.

\begin{theorem}[$*$-split closure]\label{thm:A}
Let $a$ be a Hodge $(2,2)$ character of $X^4_m$, $m$ arbitrary, whose $6$-multiset splits as two zero-sum
triples $\beta'\uplus\gamma'$. Then the rational Hodge classes of $M(a)$ are algebraic:
each is the image, under the algebraic joining-line correspondence of the Shioda--Katsura
structure map, of a $\Q$-linear combination of divisor classes on $X^1_m\times X^1_m$
furnished by Lefschetz $(1,1)$ (no single ruled-surface representative is claimed).
\end{theorem}
In particular the classes at $m=39,117,195$ (two orbits each), none quasi-decomposable or
standard, are closed; the orbit of $(1,7,16,22,34,37)$ is a certified gap family ($\nu\notin S_m$ at $39$ and at its
inflations $117,195$ --- all verifier-backed), not covered by
the published mechanisms compared in \S\ref{sec:dossier} (to the author's knowledge), while the
second orbit's $\nu$ lies in $S_{39}$ (an explicit six-generator $\pm1$ certificate), so its algebraicity is
also derivable from the standard calculus (Lemma~4.1(iii) of \cite{Aoki00}) --- the closure certificates and the join-transport construction scheme are new for both.

\begin{theorem}[two-pair standard closure]\label{thm:Aprime}
Let $a$ be a Hodge $(2,2)$ character of $X^4_m$ admitting vanishing pairs $p_1=\{k,m-k\}$,
$p_2=\{\ell,m-\ell\}$ with
\[
a\uplus p_1\uplus p_2=S\uplus Q\qquad(\text{equality of }10\text{-element multisets}),
\]
where $S$ is an Aoki $5$-standard sextuple and $Q$ a grade-$2$ Hodge quadruple. Then the rational Hodge
classes of $M(a)$ are algebraic.
\end{theorem}
This applies to---and, by exhaustive search, uniquely presents---the three remaining $5\mid m$ classes of
the census:
\begin{center}\small
\begin{tabular}{llll}
\toprule
$m$ & $a$ & $p_1,p_2$ & $S\uplus Q$\\
\midrule
$45$ & $(1,19,20,28,30,37)$ & $(5,40),(10,35)$ & $(1,10,19,28,37,40)\uplus(5,20,30,35)$\\
$105$ & $(3,24,50,66,85,87)$ & $(15,90),(45,60)$ & $(3,24,45,66,87,90)\uplus(15,50,60,85)$\\
$105$ & $(1,22,43,64,90,95)$ & $(5,100),(20,85)$ & $(1,22,43,64,85,100)\uplus(5,20,90,95)$\\
\bottomrule
\end{tabular}
\end{center}
---none quasi-decomposable, standard, or Theorem~\ref{thm:A}-splittable: closing, with the
induced $m=135$ copy, the $p=5$-standard-block classes of the census. Of these, the
$\Q(\zt_{21})$ orbit at $105$ is a certified gap class, not covered by the published mechanisms
compared in \S\ref{sec:dossier} (to the author's knowledge); $m=45$ and
$135$ fall under the degree forms of \cite[Thm.~0.1(i)]{Aoki00}, and the $\Q(\zt_7)$ orbit's $\nu$
lies in $S_{105}$, so for those three the algebraicity is also derivable from \cite{Aoki00} and the
explicit two-pair presentations are the new content.

\begin{theorem}[level-lifted closure; the $m=33$ witness]\label{thm:App}
Let $w=(1,4,16,22,25,31)$, the unique non-quasi-decomposable non-standard Hodge class of $X^4_{33}$
(uniqueness is the census's, used for Theorems~\ref{thm:B} and~\ref{thm:C}; the proof below uses
only the displayed identities),
and $a=2w\bmod 66$ its induced copy on $X^4_{66}$. Then
\[
a\uplus\{1,65\}\uplus\{25,41\}=Q\uplus S,\qquad Q=(1,25,44,62),\ S=(2,8,32,41,50,65),
\]
with $Q$ a grade-$2$ Hodge quadruple and $S$ quasi-decomposable at level $66$ via
$S\uplus\{33,33\}=(2,32,33,65)\uplus(8,33,41,50)$, the self-paired element $(33,33)\in M_{66}(1)$
being its unique quasi-witness. Hence $\V(a)$ is algebraic (Lefschetz $+$ da Silva $+$ Aoki~1-4, as in
Theorem~\ref{thm:Aprime}), and $\V(w)$ is algebraic by descent along $\pi\colon X^4_{66}\to X^4_{33}$,
$(x_i)\mapsto(x_i^2)$ (Lemma~\ref{lem:descent}).
\end{theorem}

\begin{theorem}[completion of the odd-degree census; $m\le199$]\label{thm:B}
For every odd $m\le199$, every grade-$3$ Hodge $(2,2)$ class of $X^4_m$ is algebraic: the classes
beyond $\text{decomposable}\cup\text{quasi}\cup\text{standard}$ are exhausted by
Theorems~\ref{thm:A}, \ref{thm:Aprime}, \ref{thm:App} and level transport (the inflation/descent
pair). In particular the Hodge conjecture holds \emph{in full} for every Fermat fourfold of
odd degree $\le199$: in codimension $2$ the non-primitive summand of $H^4$ is spanned by the
algebraic class $h^2$, so the statement follows from the primitive one; codimension $1$ is
Lefschetz $(1,1)$; and codimension $3$ follows from codimension $1$ by hard Lefschetz ---
$L^2\colon H^2(\Q)\xrightarrow{\ \sim\ }H^6(\Q)$ is an isomorphism of Hodge structures, so
every $(3,3)$ Hodge class is $L^2$ of a $(1,1)$ Hodge class, and $L^2$ of a divisor class is
algebraic. (Codimensions $0$ and $4$ are trivial: $H^0$ and $H^8$ are spanned by the classes of
$X$ and of a point.)
\end{theorem}

\begin{theorem}[field obstruction; Question~1]\label{thm:C}
Let $a_0=(7,10,13,19,22,28)=7w$ be da Silva's representative of the $m=33$ orbit. No algebraic cycle
defined over $\Q(\zt_{33})$ has nonzero projection onto the Galois orbit of $\V(a_0)$: Frobenius at any
prime above $67$ acts on every conjugate eigenline by a primitive sixth root of unity, of exact
order $6$ --- an exact certificate $u(a_0,67)=1+\zt_3=\zt_6$ (for the pinned character
choice of \S\ref{sec:C}) computed
in $\Z[x]/\Phi_{33}$ with no floating point. (Under an alternative Jacobi-sum definition
$j_{\mathrm{alt}}=-j$ the formula relating $j_{\mathrm{alt}}$ to geometric Frobenius acquires
the compensating sign; the geometric eigenvalue remains $\zt_6^{\pm1}$.) Since $W$ has $\Q$-coefficients, $[W]$ and all its translates
have zero projection: \emph{the answer to Question~1 is NO at $m=33$}.
\end{theorem}
Consequently the obstruction is local and exact: if a cycle defined over a finite extension
$L/\Q(\zt_{33})$ --- not assumed Galois over $\Q(\zt_{33})$; the constraint is per prime ---
has nonzero projection, then at every prime
$\mathfrak q\mid\mathfrak p\mid67$ of $L$ one needs
$u(a_0,67)^{f(\mathfrak q/\mathfrak p)}=1$ --- no hypothesis on the ramification of
$L/\Q(\zt_{33})$ is needed: the eigenline's character is unramified at $67$ over the base, so
inertia at $\mathfrak q$ acts trivially and Frobenius acts through
$u^{f(\mathfrak q/\mathfrak p)}$ --- so $6\mid f(\mathfrak q/\mathfrak p)$: every
residue degree of $L$ above $67$ is divisible by $6$. (The same machinery, run at the closed
levels, computes values of order $9$ at $m=45$ and of orders $42$ and $7$ at $m=105$ at sampled
split primes --- computational observations, not part of any theorem here; they indicate that
the cycles of Theorem~\ref{thm:Aprime} likewise cannot be $\Q(\zt_m)$-rational.)

\begin{proposition}[standard $=$ AP]\label{prop:D}
Let $m$ be odd. Aoki's standard characters $\sigma_{p,x}$ for an odd prime $p\mid m$ --- defined for
$x$ with $d/\gcd(x,d)>2$, $d=m/p$ \cite[\S1]{Aoki87} --- are exactly the padded arithmetic-progression
characters $\{x,x+m/p,\dots,x+(p-1)m/p,-px\}$ with $px\not\equiv0\pmod m$, of grade $(p+1)/2$;
grade-$3$ standard classes exist iff $5\mid m$ and $m>5$.
\end{proposition}

\medskip
\noindent\emph{Verification status.} Theorem~\ref{thm:A}: proved (transport pinned to the Shioda--Katsura
primary text). Theorem~\ref{thm:Aprime}: proved (a citation chain through Aoki's
Theorems~1-1, 1-4(i),(ii), 2-1 \cite{Aoki87} $+$ Lefschetz $(1,1)$ $+$ an inflation lemma; all
identities verified exactly, including $u(a,p)=u(S,p)u(Q,p)$).
Theorem~\ref{thm:B}: machine-verified census, reproduced exactly by an algorithmically
independent implementation (the four-tier certificate structure is laid out in \S\ref{sec:B}).
Theorem~\ref{thm:C}: proved, exact certificate. Proposition~\ref{prop:D}:
read off the primary construction. We give no explicit defining equations for a cycle on $X^4_{33}$
(Theorem~\ref{thm:App}'s certifying cycles are constructed at level $66$ and transported by the finite
morphism).

\section{Standard cycles are the AP family (Proposition~\ref{prop:D})}\label{sec:D}
For odd $p\mid m$, $d=m/p$, and $x$ with $d/\gcd(x,d)>2$, the standard character is the multiset
$\{x,x+d,\dots,x+(p-1)d,\,m-px\}$ of grade $(p+1)/2$ (Aoki \cite[\S1]{Aoki87}; da Silva \cite{daSilva},
Appendix). For \emph{odd} $m$ the admissibility hypothesis is equivalent to the visible one:
$d$ is odd, so $d/\gcd(x,d)\le2$ forces $d\mid x$, i.e.\ $px\equiv0\pmod m$ --- exactly the degenerate
case excluded in the statement. Grade $3$ forces $p=5$; at $m=5$ every candidate degenerates
($5x\equiv0$), whence ``$m>5$''. The identification is load-bearing for
Theorem~\ref{thm:B}: it is what makes the census's standard classifier complete for odd $m$
(Lemma~\ref{lem:complete}).
The padded arithmetic-progression family is exactly the witness family of the field-of-definition
threshold of the companion papers: at the surface level in \cite{P1} (with explicit Hasse--Davenport
witnesses), and in general even dimension in \emph{Galois-invariant ranks of middle algebraic cycles
on Fermat varieties} (in preparation). \qed

\begin{remark}[even $m$]\label{rem:evenstd}
For even $m$ (used only in the even census of \S\ref{sec:even}) the admissibility hypothesis is
strictly stronger than $px\not\equiv0$: it excludes $x\equiv d/2\pmod d$. Every excluded multiset
contains a vanishing pair --- $(x+k_1d)+(x+k_2d)\equiv0\pmod m$ reduces to $1+2j+k_1+k_2\equiv0\pmod5$
for $x=d/2+jd$, solvable with $0\le k_1<k_2\le4$ since $k_1+k_2$ ranges over all residues mod~$5$ ---
so it is decomposable, hence algebraic for the trivial reason, and the census classifier (which tests
decomposability first) is unaffected. (Only $p=5$ yields grade-$3$ sextuples, so this is the only
case the classifier meets.)
\end{remark}

\section{The \texorpdfstring{$*$}{*}-split closure theorem (Theorem~\ref{thm:A})}\label{sec:A}

The proof factors through the following transport lemma, which records separately the passage from complex eigenlines to rational algebraic classes: which objects
are correspondences, over which fields they are defined, where the Tate twist sits, and why the
scalar onto $\V(a)$ is nonzero.

\begin{lemma}[rational transport]\label{lem:transport}
Let $C=X^1_m$ ($m$ arbitrary), $a=\beta'\!*\gamma'$ a Hodge $(2,2)$ character split into two zero-sum
triples, and let $f$ denote the second summand of the Shioda--Katsura structure map
\cite[Thm.~2.2]{daSilva} (going back to Shioda's inductive structure \cite{Shioda79}; the
geometric join construction is \cite[Thm.~1.7]{SK}), a $G$-equivariant morphism of Hodge structures
\[
f\colon H^1_{\mathrm{prim}}(C,\Q)\otimes H^1_{\mathrm{prim}}(C,\Q)\;\longrightarrow\;
H^4_{\mathrm{prim}}(X^4_m,\Q)
\]
of type $(1,1)$ --- equivalently, a weight-preserving $(0,0)$-morphism out of
$H^1\otimes H^1(-1)$; this $(-1)$ is the only Tate twist in the construction. Then:
\begin{enumerate}
\item[(i)] at every $t\in(\Z/m)^\times$ the line $\V(t\beta')\otimes\V(t\gamma')$ has Hodge type
$(1,0)\otimes(0,1)$ or $(0,1)\otimes(1,0)$, so the block
$B(a):=\bigoplus\V(t\beta')\otimes\V(t\gamma')\subset H^2(C\times C,\C)$ --- the sum over the
distinct \emph{ordered} paired characters $(t\beta',t\gamma')$, pairs of ordered triples, whose
coordinatewise stabilizer may be nontrivial (e.g.\ $\beta'=(3,3,3)$ at $m=9$) and coincides with
the ordered stabilizer of $a=\beta'*\gamma'$, so the summands of $B(a)$ and $M(a)$ biject ---
is purely of type
$(1,1)$;
\item[(ii)] $B(a)$ is Galois-stable, so $B(a)_\Q:=B(a)\cap H^2(C\times C,\Q)$ satisfies
$B(a)_\Q\otimes\C=B(a)$, and by Lefschetz $(1,1)$ every element of $B(a)_\Q$ is a $\Q$-combination
of divisor classes on the surface $C\times C$;
\item[(iii)] $f$ is induced by an algebraic correspondence (the blow-up--quotient--blow-down chain
of \cite{SK}, defined over the cyclotomic field of the join construction), and it is a
$G$-equivariant \emph{isomorphism} onto its structural summand \cite[Thm.~2.2]{daSilva}: it
maps each one-dimensional eigenline $\V(t\beta')\otimes\V(t\gamma')$ isomorphically --- in
particular by a nonzero scalar --- onto $\V(ta)$, hence carries $B(a)$ onto the corresponding
block of $M(a)$. On decomposable \emph{algebraic} inputs $f(Z_1\otimes Z_2)=m\,Z_1\wedge Z_2$,
the ruled join by lines \cite[Thm.~2.2(c)]{daSilva}; this formula is used only for the geometric
description of the image cycles, not for the scalar on the $H^1$-eigenlines, which are not
spanned by algebraic classes;
\item[(iv)] consequently every rational Hodge class of $M(a)$ is the class of an algebraic cycle:
it is the $f$-image of a $\Q$-combination of divisor classes --- the image of an algebraic
class under the algebraic joining-line correspondence, hence an algebraic cycle class over
$\Qbar$, lying in $H^4(X^4_m,\Q)$ by construction.
Field-of-definition bookkeeping: the correspondence above is defined over a cyclotomic field, and
no $\Q$-rationality of the \emph{cycles} is claimed or needed --- algebraicity of a rational Hodge
class is a $\Qbar$-statement (the fields the certifying cycles live over are the subject of
Theorem~\ref{thm:C}, not of this lemma).
\end{enumerate}
\end{lemma}

\begin{proof}
(i) is the grade bookkeeping $|t\beta'|+|t\gamma'|=3$, each summand in $\{1,2\}$ (verified exactly
for every printed instance in the ancillary identities). (ii): Galois permutes the summands of
$B(a)$ (the index set is one orbit), so $B(a)$ is defined over $\Q$ as a sub-Hodge structure; being
purely $(1,1)$, its rational part is divisorial by Lefschetz $(1,1)$. (iii) is \cite[Thm.~2.2]{daSilva} (cf.\ \cite{Shioda79}, \cite[Thm.~1.7]{SK}): $f$ is a $G$-equivariant
isomorphism onto its image, so its restriction to each one-dimensional eigenline is injective,
i.e.\ a nonzero scalar, and equivariance matches the source and target characters. (iv)
combines (i)--(iii). (An explicit block projector --- an orbit sum of graph correspondences
with rational, Galois-symmetric coefficients, algebraic over $\Q$ --- can be written down, but
is not needed: the class handed to $f$ already lies in $B(a)_\Q$.)
\end{proof}

\begin{proof}[Proof of Theorem~\ref{thm:A}]
A split $a=\beta'\uplus\gamma'$ into zero-sum triples realizes $a\sim\beta'*\gamma'$, and by
multiplicity one $M(a)$ is the $f$-image of the block $B(a)$ of Lemma~\ref{lem:transport}. By the
lemma, every rational Hodge class of $M(a)$ is the $f$-image of a $\Q$-combination of divisor
classes on $C\times C$, with a nonzero scalar on each eigenline (part (iii) of the lemma) ---
algebraic, because the image of an algebraic class under an algebraic correspondence is
algebraic. Geometrically the correspondence is the joining-line construction --- over
$(z,w)\in C\times C$ the line $\lambda z+\mu w$ lies on $X^4_m$:
$\sum(\lambda z_i)^m+\sum(\mu w_j)^m=\lambda^m\cdot0+\mu^m\cdot0$ --- and on decomposable
algebraic inputs its effect is the ruled join $f(Z_1\otimes Z_2)=m\,Z_1\wedge Z_2$
\cite[Thm.~2.2(c)]{daSilva}; for a general divisor no single ruled-surface representative of
the image class is claimed, and only the correspondence-image statement is used.
\end{proof}
\begin{remark}
Quasi-decomposability tests only decompositions with both factors Hodge; curves carry no Hodge classes in
$H^1$, so the $(P1)/(P2)$ formalism is structurally blind to this odd$\otimes$odd branch. At $m=39$:
$(1,7,16,22,34,37)=(1,16,22)\uplus(7,34,37)$ and $(1,14,16,22,29,35)=(1,16,22)\uplus(14,29,35)$; the
second has $j(a)=p^2$ exactly at every computed split prime --- consistent with a divisorial
certificate whose class is fixed by the Frobenius elements in question (algebraicity alone does
not force this: a divisor defined only over an extension can carry a nontrivial finite
character), the calibration promised in the Notation.
\end{remark}

\section{The two-pair closure theorem (Theorem~\ref{thm:Aprime})}\label{sec:Aprime}
\begin{proof}
Write $\delta=p_1*p_2=(k,m-k,\ell,m-\ell)\in\Ds^2_m$ (Aoki's all-pairs classes, represented by linear
cycles with nonzero self-pairing, \cite[Thm.~1-1]{Aoki87}). The multiset identity gives $a*\delta\sim S*Q$
on $X^8_m$ ($\sim$ a coordinate permutation, an algebraic automorphism).
(1)~$Q$ is a grade-$2$ Hodge quadruple: its rational Galois block on the Fermat surface is
purely of type $(1,1)$, hence divisorial by Lefschetz $(1,1)$, and claim$(Q)$ follows.
(2)~For the $\sigma_{5,1}$ blocks ($m=45,d=9$ and $m=105,d=21$; $\gcd(1,d)=1$) claim$(S)$ is
\cite[Thm.~2-1]{Aoki87}; for the $\sigma_{5,3}$ block at $m=105$ note $S=3\cdot\sigma_{5,1}$ at level $35$
and apply Lemma~\ref{lem:infl}. (In general a $5$-standard block $\sigma_{5,x}$ with
$g=\gcd(x,d)>1$ is the $g$-inflation of $\sigma_{5,x/g}$ at level $m/g$, where the gcd is $1$, so
Lemma~\ref{lem:infl} and Thm.~2-1 cover every admissible block.)
(3)~claim$(S*Q)$ is \cite[Thm.~1-4(i)]{Aoki87} ($r=4$, $s=2$, $n=8$).
(4)~claim$(a)$ is \cite[Thm.~1-4(ii)]{Aoki87}: \emph{if there exists $\delta\in\Ds^s_m$ with
claim$(\alpha*\delta)$, then claim$(\alpha)$}---applied with $\alpha=a$, $s=2$.
\end{proof}
\begin{lemma}[inflation]\label{lem:infl}
Let $m=e\,m'$ and $\beta$ a level-$m'$ character with every entry of $e\beta$ nonzero mod $m$. Then
$\mathrm{claim}_{m'}(\beta)\Rightarrow\mathrm{claim}_m(e\beta)$.
\end{lemma}
\begin{proof}
$\pi\colon X^n_m\to X^n_{m'}$, $(x_i)\mapsto(x_i^e)$ is a morphism with $\pi(\zt x)=\zt^e\pi(x)$ for
$\zt\in\mu_m$, so $\pi^*$ maps $\V(\beta)$ into $\V(e\beta)$ equivariantly (the nonzero-entries
hypothesis is what makes $e\beta$ a legitimate primitive character, so the target eigenline
exists); $\pi$ finite surjective gives
$\pi_*\pi^*=(\deg\pi)\,\mathrm{id}\neq0$ on $\Q$-cohomology, so $\pi^*\xi$ spans $\V(e\beta)$, and for a
representing cycle $Z$, $\omega_{e\beta}(\pi^*Z)=\pi^*(\omega_\beta(Z))\neq0$. (Here $\pi^*Z$
is the pullback in $CH_\Q$: $\pi$ is a finite surjective morphism of smooth projective
varieties, hence flat by miracle flatness, and flat pullback is compatible with cycle classes.)
\end{proof}
\begin{remark}
da Silva's quasi-decomposability adjoins one pair and splits $4+4$; the two-pair $6+4$ pattern is outside
the definition. The three presentations were found by exhaustive iterated-closure search and are unique
per class; equivalently the three classes are sum-preserving two-entry surgeries of standard sextuples
($(10,40)\mapsto(20,30)$ at $45$; $(45,90)\mapsto(50,85)$, $(85,100)\mapsto(90,95)$ at $105$). The characters multiply
exactly: $\chi_a=\chi_S\chi_Q$, i.e.\ $u(a,p)=u(S,p)\,u(Q,p)$ at every split prime. The surviving $m=33$ class
resists this genre exhaustively (no augmentation by up to five vanishing pairs admits a certifying
partition; $5\nmid33$, so no standard blocks exist).
\end{remark}

\section{Level-lifted closure of the witness (Theorem~\ref{thm:App})}\label{sec:App}
\begin{proof}
All identities are exact (machine-verified; the ancillary re-derives each). The augmented multiset
partitions as stated; $Q$ is grade-$2$ Hodge. For $S$: the pair $(33,33)$ is a legitimate element of
$M_{66}(1)$ ($\fr{33t}+\fr{33t}=66$ for every odd $t$; multiplicities are permitted), and in da
Silva's coordinates $S=\beta\#\gamma$ with $\beta=(2,32,65;33)$,
$\gamma=(8,41,50;33)$ --- an element of his $U^{2,2}_{66}$, the distinguished entries $33,33$ summing
to $0\bmod66$ (the definition preceding \cite[Cor.~2.3]{daSilva}); Corollary~2.3(b) with $(r,s)=(2,2)$
and Lefschetz~$(1,1)$ for $\beta,\gamma$ give claim$(S)$. The hypotheses of Corollary~2.3(b) are
exactly: $r,s$ even, algebraicity of both factors, nonzero entries, and the matching condition ---
each verified above; the self-paired value $b_3=c_3=33=m/2$ is admissible
(Remark~\ref{rem:selfpair}). Then, exactly as in
Theorem~\ref{thm:Aprime}: claim$(S*Q)$ by \cite[Thm.~1-4(i)]{Aoki87}; $S*Q\sim a*\delta$ with
$\delta=(1,65,25,41)\in\Ds^2_{66}$; claim$(a)$ by \cite[Thm.~1-4(ii)]{Aoki87}. Descent is
Lemma~\ref{lem:descent}.
\end{proof}
\begin{remark}[self-pair admissibility]\label{rem:selfpair}
The hypotheses of the pairing statement used above --- membership in $U^{2,2}_{66}$ and
Corollary~2.3(b) of \cite{daSilva} --- are: even $r,s$, nonzero entries, algebraicity of both
factors, and the matching condition $b_3+c_3\equiv0$. Pairwise distinctness of the entries is
not among them, and the underlying structure isomorphism \cite[Thm.~2.2]{daSilva} carries no
nondegeneracy hypothesis on the characters; the self-paired distinguished value
$b_3=c_3=33=m/2$ is therefore legitimate. This is the one point where the even level is
essential: no self-pair exists at an odd level (Remark~\ref{rem:App}(iv)).
\end{remark}
\begin{lemma}[descent]\label{lem:descent}
Let $m=e\,m'$ and $\beta$ a level-$m'$ character with every entry of $e\beta$ nonzero mod $m$. Then
$\mathrm{claim}_m(e\beta)\Rightarrow\mathrm{claim}_{m'}(\beta)$.
\end{lemma}
\begin{proof}
$\pi\colon X^n_m\to X^n_{m'}$, $(x_i)\mapsto(x_i^e)$ is a finite surjective morphism
($\sum x_i^m=\sum(x_i^e)^{m'}$ identically) intertwining the group actions along the $e$-th-power
surjection $G_m\twoheadrightarrow G_{m'}$; hence $\pi_*\V(e\beta)\subseteq\V(\beta)$ (again the
nonzero-entries hypothesis makes $e\beta$ a legitimate primitive character). Since
$\pi_*\pi^*=(\deg\pi)\,\mathrm{id}\neq0$ and $\pi^*\V(\beta)=\V(e\beta)$ with both eigenlines
$1$-dimensional, $\pi_*\colon\V(e\beta)\xrightarrow{\ \sim\ }\V(\beta)$. Proper pushforward carries
algebraic cycle classes to algebraic cycle classes ($\mathrm{cl}\circ\pi_*=\pi_*\circ\mathrm{cl}$),
and $\pi$ is finite of relative dimension zero, so there is no degree shift or Tate twist.
\end{proof}
\begin{remark}\label{rem:App}
(i) With the inflation lemma (Lemma~\ref{lem:infl}), claim is level-transparent:
$\mathrm{claim}_{m'}(\beta)\Leftrightarrow\mathrm{claim}_m(e\beta)$ --- every lift of an open class
is an independent chance to close it. (ii) The mechanism is selective: the two-pair-at-the-double search does not fire for the induced
copies of the $m=54$ classes at $108$ or of the $m=50$ class at $100$ (all of these close at their own
levels at depth $3$--$4$; even-census receipts, \S\ref{sec:even}), and the beyond-machinery $m=66$ classes resist the
single two-pair at level $66$ and its lift $132$ (they too close at depth $3$;
even-census receipts, \S\ref{sec:even}). (iii) Consistency with
the field obstruction (\S\ref{sec:C}): the constructed cycles (Lefschetz divisors on $X^2_{66}$,
Aoki's linear spaces with $\varepsilon=\zt_{132}$, radicals) are not $\Q(\zt_{66})$-rational, and
$\Q(\zt_{66})=\Q(\zt_{33})$ --- the obstruction forbids exactly the $\Q(\zt_{33})$-rational
certificates, and this theorem produces non-rational ones (their precise minimal field of
definition is not computed here). The local Frobenius values agree exactly:
$u(2w@66,67)=\zt_6=u(w@33,67)$. (iv) The
uniqueness of the quasi-witness $v=33$ shows the self-paired $m/2$ element --- which exists only at
even levels --- is strictly load-bearing: this is why the exhaustive level-$33$ search (five
augmentation pairs, all partition types) could not see the closure.
\end{remark}

\section{The census and its completeness (Theorem~\ref{thm:B})}\label{sec:B}
For each odd $m$: enumerate grade-$3$ Hodge multisets (meet-in-the-middle on half-unit sum profiles);
reduce to Galois orbits; classify decomposable $\to$ quasi $\to$ standard (Prop.~\ref{prop:D}) $\to$
$*$-split (Thm.~\ref{thm:A}) $\to$ survivors; the surviving orbits are then closed by
Theorems~\ref{thm:Aprime}/\ref{thm:App} and level transport. The engine reproduces
the Shioda/da Silva anchors ($(P_m)$ at $m=9,21,27$; the $m=33$ witness) and was reproduced, on all $89$
odd levels of the census list ($21\le m\le199$, $m\neq23$; the omitted small levels and $m=23$ are
classical \cite[Thms.~2.6--2.7]{daSilva}), by a second, algorithmically independent
implementation (structured-array meet-in-the-middle,
cross-validated against its own brute force for $m\le143$). After Theorem~\ref{thm:A} (which closes the
$*$-split classes at $39,117,195$), seven beyond-machinery orbits remain: four primitive
($m=33,45,105\times2$) and the induced copies at $99,135,165$; after Theorem~\ref{thm:Aprime} (which closes $45$ and both $105$, hence by Lemma~\ref{lem:infl}
also $m=135$) and Theorem~\ref{thm:App} (which closes $33$, hence $99$ and $165$),
nothing remains. Induced copies are tied to their primitives by the finite pull--push
correspondences of Lemmas~\ref{lem:infl} and~\ref{lem:descent}, which transfer claim in both
directions.

\begin{lemma}[census completeness]\label{lem:complete}
Fix $m$ and consider the shipped generator, classifiers, and canonizer.
\begin{enumerate}
\item[(i)] Every element of $M_m$ has even length $2g$, $g$ its grade: for nonzero entries
$\fr{(m-t)x}=m-\fr{tx}$, so $\sum_i\fr{(m-t)a_i}=Lm-\sum_i\fr{ta_i}$ and grade-constancy forces
$g=L/2$. In particular $M_m(1)$ consists exactly of the vanishing pairs $\{k,m-k\}$ (including the
self-pair $k=m/2$ at even $m$), and a grade-$3$ sextuple is decomposable exactly when it
contains a vanishing pair --- the complementary quadruple is then automatically grade-$2$ Hodge.
\item[(ii)] The quasi-decomposability test is complete for da Silva's definition: for a pair-free
sextuple $a$, the $8$-multiset $a\uplus\{k,m-k\}$ splits in $M_m$ only with part-lengths $(2,6)$ or
$(4,4)$, by (i); in every $(2,6)$ split the pair is either $\{k,m-k\}$ itself or $\{a_i,m-a_i\}$
with $k\in\{a_i,m-a_i\}$, and in both cases the split is the trivial one (the pair part is the
adjoined pair, the sextuple part is $a$ itself) --- excluded by the non-triviality clause of
\cite[Def.~2.4]{daSilva}. Hence $a$ is quasi-decomposable iff some
$a\uplus\{k,m-k\}$ splits into two grade-$2$ Hodge quadruples --- exactly the test performed.
\item[(iii)] The meet-in-the-middle generator produces exactly the grade-$3$ Hodge sextuples: the
join condition enforces $\sum_i\fr{ta_i}=3m$ for every $t$ in a half-unit system, and the
complement identity of (i) extends it to all units; conversely every sorted sextuple is the join of
its two sorted halves, each of which occurs in the enumeration of \emph{all} $3$-multisets of
nonzero residues.
\item[(iv)] Canonization retains the lexicographically least sorted conjugate --- a system of
distinct representatives for the $(\Z/m)^\times$-orbits.
\end{enumerate}
Consequently the census verdicts (decomposable / quasi / standard / $*$-split / survivor) coincide
with their definitions level by level; the standard classifier is complete for odd $m$ by
Proposition~\ref{prop:D}.
\end{lemma}

\begin{proof}
(i) is the two displayed identities; for the last claim, a vanishing pair contributes $m$ to
every unit sum, so the complementary quadruple has constant grade $2$. (ii): the parts have even lengths summing to $8$ with grades
$(1,3)$ or $(2,2)$; in a $(2,6)$ split the length-$2$ part is a vanishing pair inside
$a\uplus\{k,m-k\}$, and since $a$ is pair-free it is $\{k,m-k\}$ (complement $=a$) or uses one
adjoined entry, forcing $k\in\{a_i,m-a_i\}$ and complement $(a\smallsetminus\{a_i\})\uplus\{a_i\}=a$
(multiplicities included; the even-$m$ self-pair case is identical). (iii): if the joined profile
sums to $3m$ on the half-units $t$, then at $m-t$ it sums to $6m-3m=3m$; the converse containment
is the split of a sorted sextuple at position $3$. (iv): distinct orbits have distinct minima.
The implementations are additionally calibrated against direct brute force (the independent
reimplementation for $m\le143$; the even-parity engine, exact match of representative counts,
for $m\le14$) and reproduce the Shioda/da~Silva anchors above. The full $89$-level run ---
per-level counts, classification tallies, survivor lists, and per-orbit closure witnesses ---
ships as checksummed data artifacts with the ancillary files.
\end{proof}

\emph{What is certified, and how.} Four tiers, in decreasing logical strength. (1)~\emph{Per-orbit
certificates}: for each of the $78{,}299$ orbit representatives across the $89$ levels, a stored
witness of its classification (the vanishing pair / the quasi split / the standard parameter /
the $*$-split triples), re-verified against the definitions on every run of the supplied runner;
for the seven beyond-machinery orbits the runner re-establishes the \emph{negative} screenings
live (no vanishing pair, no quasi witness at any $k$, not standard, no zero-sum split), so the
terminal labels are re-checked certificates, not labels. (2)~The \emph{completeness lemma}
above, which reduces ``no orbit is missed'' to the definitional correctness of the shipped
generator and canonizer. (3)~The \emph{independent reimplementation}: an algorithmically
independent second program regenerates the census from scratch and matches it level by level
(its complete $89$-level log ships; the runner re-executes it live at anchor levels and
regenerates a fresh sub-range against the summary table below). (4)~\emph{Checksummed receipts}
of the historical campaign runs: a checksum certifies file integrity, never recomputation ---
recomputation is what tiers (1)--(3) execute. A per-level summary table
(\texttt{census\_level\_summaries.json}: representative count, tally by kind, survivor list,
and the SHA-256 hash of the canonically sorted representative list) makes ``no missed orbit''
checkable level by level: any regeneration, by any implementation, must reproduce the count and
the hash.

\emph{The bound $199$.} The bound is a computational horizon, not a structural one: every
closure mechanism used here is degree-uniform, and nothing in the method stops at $199$. What
grows is the census cost --- on the reference machine (Apple M4 Max, Python 3.13) the engine
takes $158$\,s at the single level $m=199$, the full two-engine $89$-level sweep of the pinned
receipt took $\approx36$ minutes, and the complete independent rerun $\approx12$ minutes --- and $199$ was
chosen to include, with margin, the largest levels at which new primitive phenomena appear in
the odd census: the two-pair genre at $105$ (the last new primitive beyond-machinery orbit) and
the $*$-split gap family through $195=5\cdot39$. The probes beyond the bound (next paragraph)
found no new phenomenon.

Additional
single-engine censuses at the induced-risk levels beyond $199$ --- computational observations:
these three levels sit outside the witness-tier receipts of Theorem~\ref{thm:B} --- found: at
$231=33\cdot7$ and
$297=33\cdot9$ one beyond-machinery orbit each, the induced copy of the $m=33$ class (content $7$,
resp.\ $9$), closed by Theorem~\ref{thm:App} and transport; at the control level $273=3\cdot7\cdot13$,
two $*$-split classes (closed by Theorem~\ref{thm:A}) and nothing else. They indicate that the
conclusion of Theorem~\ref{thm:B} also holds at $m=231,273,297$; the theorem itself does not
depend on these levels.

\section{Computational outlook: the even-degree boundary through degree 250}\label{sec:even}

The odd-degree theorems are complete. The even-degree boundary is a different landscape, and this
closing section only \emph{reports its computed map} --- engines, closure ledger, certificates,
and pinned receipts ship in the ancillary --- deferring systematic treatment (full statements and
proofs of the closure and wall structure) to the companion paper \cite{BNeven}, with which this
note shares its ancillary package. Nothing in Theorems~\ref{thm:A}--\ref{thm:C} or in the odd
census depends on this section, and every computational claim below is backed
receipt-by-receipt within the ancillary of \emph{this} submission --- the companion adds
statements and proofs, not data. All censuses use the
parity-corrected decomposability test (the self-paired element $m/2$; part~(i) of
Lemma~\ref{lem:complete} covers both parities); every lattice verdict below carries an
independent certificate in the ancillary --- a modular kernel witness for each non-membership, a
re-summed integer combination for each membership. One statement is load-bearing also for the
odd-degree novelty classification above, and is therefore proved here.

\begin{proposition}[the implemented lattice is Aoki's $S_m$]\label{prop:ident}
Aoki \cite[\S2]{Aoki00} defines $A_m$ as the zero-sum sublattice of the free abelian group on
$\Z/m\smallsetminus\{0\}$; $u$ --- our $\nu$ --- as the multiplicity-vector map, additive
under juxtaposition, $\nu(\alpha*\beta)=\nu(\alpha)+\nu(\beta)$; the standard elements,
verbatim, as
$\sigma_{p,a}=(a,\,a+\frac mp,\dots,a+\frac{(p-1)m}p,\,m-pa)$ for $p\ge3$ and
$\sigma_{2,a}=(a,\,a+\frac m2,\,m-2a,\,\frac m2)$ for $p=2$, over primes $p\mid m$ ($m>p$) and
$0<a<\frac mp$ with nonzero entries; $\mathfrak S_m$ as the set of $\alpha$ admitting
$\alpha*\delta\sim\sigma_1*\cdots*\sigma_k*\delta'$ with $\delta,\delta'$ unions of vanishing
pairs, with $\mathfrak D_m\subset\mathfrak S_m$; and $S_m\subset A_m$ as the subgroup generated
by $\nu(\mathfrak S_m)$. Consequently $S_m$ equals the $\Z$-span of the pair vectors
$\nu(\{k,m-k\})$ and the standard vectors $\nu(\sigma_{p,a})$, $p=2$ included --- exactly the
generator list of the implemented lattice --- so the machine verdicts $\nu\in S_m$ /
$\nu\notin S_m$ are statements about Aoki's group itself.
\end{proposition}
\begin{proof}
$\supseteq$: every pair class lies in $\mathfrak D_m\subset\mathfrak S_m$, and every standard
element $\sigma$ satisfies $\sigma*\delta\sim\sigma*\delta$, so $\sigma\in\mathfrak S_m$: each
listed generator is in $\nu(\mathfrak S_m)$. $\subseteq$: if
$\alpha*\delta\sim\sigma_1*\cdots*\sigma_k*\delta'$, then
$\nu(\alpha)=\sum_i\nu(\sigma_i)+\nu(\delta')-\nu(\delta)$, an integer combination of the listed
generators. The two spans coincide.
\end{proof}

\emph{The computed map} {\rm(computed; engines, closure ledger, and receipts in the ancillary;
full statements and proofs in \cite{BNeven})}. For even $m\le108$ the census finds twenty-one
primitive beyond-machinery orbits (beginning at $m=50$; the odd law ``survivor
$\Rightarrow3\mid m$'' fails in the even sector); through degree $200$ it finds fourteen more,
and through degree $250$ exactly three more --- thirty-eight base-tier primitives in all.
Iterated partition certificates of the Theorem~\ref{thm:Aprime}/\ref{thm:App} kind close all
but ten; an exhaustive exchange scan --- claim-equivalence moves
$T=S\uplus A\leftrightarrow T'=S\uplus B$, $|A|=|B|\in\{2,3\}$, with exchanger
$q=A\uplus(-B)$ a grade-$2$ Hodge quadruple (Lefschetz, unconditional) or a grade-$3$ sextuple
of known claim, \emph{every} move out of \emph{every} open class enumerated at its own level
and its lifts --- organizes the residual into claim-equivalence walls, and two further classes
then close. The $m=168$ class falls under Aoki's published Theorem~0.1(i) \cite{Aoki00}
($168=2^3\cdot3\cdot7$ has his degree form; his \S4 proof opens with the bridge quoted in the
Appendix). Citing it, we record two errata of the printed text, both verified against the
primary and neither affecting the theorem as stated: the companion CMUSP \textbf{51} (2002)
paper misassigns the statement, and the $m=28$ gap-generator entry on p.~185 is a misprint (not
a zero-sum character) --- a machine-verified \emph{possible replacement} of exactly the
required join form is $\xi_{28}=(1,9,18)*(10,21,25)$, with $\nu\notin S_{28}$,
$2\nu\in S_{28}$ (whether it is the intended generator modulo $S_{28}$ is not established; the
$168$ closure does not depend on it). The isolated degree-$210$ class $(1,79,109,121,151,169)$
closes by an explicit $40$-generator $\pm1$ lattice certificate $\nu\in S_{210}$ --- pairs
and standard elements only, converted into a literal Aoki chain by the negation closure ---
whence $\V$ is algebraic by Theorems~1-4(i),(ii) of \cite{Aoki87}; the full signed combination
ships in the ancillary and is displayed, block by block, in \cite{BNeven}. What remains
\emph{unresolved by the implemented closure family and the cited published theorems} through
even degree $250$ is eight certified gap classes --- all with $\nu\notin S_m$, $2\nu\in S_m$
(the doubling holds for every Hodge class by \cite[Thm.~2.1]{Aoki00}), all outside the degree
forms of Aoki's theorem --- organized in three claim-equivalence walls: $W_{70}$ (the champion
$(1,20,24,42,61,62)$ at $70$ and its exchange partner at $210$), $W_{110}$ (three members),
$W_{114}$ (three members). Every lattice verdict carries an independent certificate in the
ancillary (a modular kernel witness for each ``no'', the re-summed $\pm1$ combination for the
``yes''); within the scanned move family the walls are exchange-isolated, and closing any
member closes its wall. Gap-ness bounds the \emph{standard} calculus, not partition search:
the $\Q(\zt_{21})$ class at $105$ closes by Theorem~\ref{thm:Aprime} while remaining a gap
class. The wall table, the per-class dossiers, the deep negative sweeps, and the independent
even-parity verification are in \cite{BNeven}.

\section{The field obstruction and Question~1 (Theorem~\ref{thm:C})}\label{sec:C}

\emph{Normalization.} Fix a split prime $p\equiv1\bmod m$ and a multiplicative character
$\chi\colon\F_p^\times\to\Qbar^\times$ of exact order $m$ (extended by $\chi(0)=0$). Set
\[
S_0(a)=\sum_{\substack{v\in(\F_p^\times)^6\\ v_0+\dots+v_5=0}}\ \prod_{i}\chi^{a_i}(v_i),
\qquad j(a)=\frac{S_0(a)}{p-1}\in\Z[\zt_m]
\]
(the scaling $v\mapsto\lambda v$ is free since $\sum a_i\equiv0$, so $S_0$ is exactly
divisible by $p-1$; the certificate asserts this). By Weil's count of points of diagonal
hypersurfaces \cite[pp.~500--502]{Weil49}, \cite{Weil52}, \emph{geometric} Frobenius at $p$
acts on the \'etale eigenspace $\V(a)\subset H^4_{\mathrm{prim}}$ by $j(a)$, with
$|j(a)|=p^2$ in every embedding --- in our even dimension ($n=4$) the point-count sign is $+1$
--- hence on $\V(a)(2)$ by $u(a,p)=j(a)/p^2$ (arithmetic Frobenius acts by the inverse scalar;
every statement in this paper is about geometric Frobenius); the conjugate eigenlines carry
$u(ta,p)$, computed by replacing $\chi$ with $\chi^t$. Under an inhomogeneous convention
($\sum v_i=1$) the Jacobi symbol changes by a sign that the point-count formula compensates;
the certificate's asserted $(u')^3=-1$ pins the geometric order regardless. A $30$-digit
numeric cross-check against the Gauss-sum product $\prod_ig(\chi^{a_i})/p^3$ corroborates
the normalization (ancillary; corroboration only --- the certificate itself is exact).
The certificate's character is pinned: $\chi$ is the character of $\F_{67}^\times$ with
$\chi(g)=\zt_{33}$ on the smallest primitive root $g=2$ (the shipped code computes discrete
logarithms base $2$); replacing $\chi$ by a conjugate character permutes the twenty conjugate
values and can exchange $\zt_6\leftrightarrow\zt_6^{-1}$ at $t=1$ --- the
convention-independent assertions are the exact order six and the two-value distribution.

\begin{proof}
Fix a prime $\ell\neq67$ and an embedding $\iota\colon\Qbar\hookrightarrow\Qbar_\ell$;
cohomology in this proof is $\ell$-adic \'etale cohomology with $\Qbar_\ell$-coefficients
(the eigenspaces correspond under comparison), and the algebraic integers
$u(ta_0,67)\in\Z[\zt_{33}]$ are read in $\Qbar_\ell$ through $\iota$; the asserted
identities $(u')^3=-1$, $(u')^6=1$, $u'\neq1$ hold in $\Z[\zt_{33}]$, hence under every
embedding, so no choice matters. The $\ell$-adic class of a cycle over $K$ is $G_K$-fixed. Take $\mathfrak p\mid67$ in $\Q(\zt_{33})$
($67\equiv1\bmod33$ splits, good reduction). $\mu_{33}\subset\F_{67}$ makes the $G$-action
$\F_{67}$-rational, so Frobenius preserves each $\V(ta_0)$ and acts on $\V(ta_0)(2)$ by
$u(ta_0,67)$ as fixed in the Normalization above. An exact
integer dynamic program in $\Z[x]/\Phi_{33}$ gives $u(a_0,67)=1+\zt_3=\zt_6$, and the $20$ unit-indexed
values $u(ta_0,67)$, $t\in(\Z/33)^\times$ --- the Frobenius scalars on the \emph{twenty} conjugate
eigenlines of $M(a_0)$: the ordered stabilizer is trivial ($7t\equiv7$ forces $t=1$), the multiset
stabilizer $\langle4\rangle$ has order $5$, so the twenty eigenlines fall into four
coordinate-permutation classes of five; the scalar --- a symmetric function of the entries --- is
constant on each class, and the four class-values coincide in pairs --- take exactly the two
values $\zt_6^{\pm1}$, each on ten eigenlines. The certificate asserts \emph{exact order six}
on every conjugate --- $(u')^3=-1$, $(u')^6=1$, $u'\neq1$ --- and the geometric scalar is
independent of the Jacobi-sum sign convention, so each eigenvalue is a primitive sixth root of
unity, hence $\neq1$. A Frobenius-fixed vector scaled by $\neq1$ vanishes componentwise.
\end{proof}
$W$ has $\Q$-coefficients and its translates are $\Q(\zt_{33})$-rational; by the census the $m=33$ class
is the unique non-quasi non-standard class at that degree, so its projection vanishes and Question~1 is
answered negatively. The YES-branch would have proved the Hodge conjecture at these classes; the NO says
only that this candidate cannot certify it, and imposes a sharp local residue-degree constraint
on any number field of definition of a certifying cycle.

\section{The closed classes' dossiers}\label{sec:dossier}
\emph{The $m=33$ witness (closed by Theorem~\ref{thm:App})}: indicated eigenvalue field
$\Q(\sqrt{-3})$ (a computed label), all
entries $\equiv1\bmod3$; its \emph{Tate-normalized} Jacobi-sum Hecke character $\chi_a$ is of
finite order --- for a grade-constant (Hodge) character the infinity type of Weil's unnormalized
$J_a$ is the norm power $N^2$ itself, so $\chi_a=J_a/N^2$ has trivial infinity type (cf.\
\cite{Weil52}) --- and its computed values at split primes lie in $\mu_6$: kernel primes occur
($p=859$: every conjugate value exactly $1$), and order $6$ is attained ($p=67$, the prime
Theorem~\ref{thm:C} uses; only the certified per-prime obstruction is used, and no global order
is claimed). The dossier's remaining role is the acceptance test for explicit cycles
(Theorem~\ref{thm:App}'s pass by construction). A Weil-class alternative would need a precise
carrier: Weil classes on abelian \emph{fourfolds} of Weil type are now unconditionally
algebraic (Markman \cite{Markman1}, with Schoen's degeneration argument \cite{Schoen}), but
the Shioda--Katsura realization carries $\V(w)$ on CM abelian varieties of dimension $\ge10$
(forced by the faithful $\Q(\zt_{33})$-action), and no algebraic compression onto a
$\Q(\sqrt{-3})$-Weil fourfold is known --- the route stays conditional on that open
compression, and is now unnecessary.

\emph{The closed classes $m=45,105$}: their exact eigenvalue fields $\Q(\zt_9),\Q(\zt_7),\Q(\zt_{21})$
(degrees $6,6,12$) exceed the imaginary-quadratic multiplication field required by the published
algebraicity theorems for Weil-type classes known to us (Schoen \cite{Schoen}; Markman
\cite{Markman1}); the CM-field secant-sheaf program \cite{Markman2}, which would reach larger fields,
is explicitly conditional on a semiregularity ``not addressed yet'' (loc.\ cit.); on
the canonical CM carriers (dimensions $12$ and $24$) restriction of scalars to an imaginary quadratic
subfield shifts the cohomological degree away from $4$, so no descent to the covered case exists.
For the $\Q(\zt_{21})$ class at $105$, Theorem~\ref{thm:Aprime}'s Fermat-internal route is,
among the mechanisms compared here, at present
the only proof of algebraicity known to the author (a gap class: $\nu\notin S_{105}$,
$2\nu\in S_{105}$); the
$m=45$ and $\Q(\zt_7)$ $m=105$ classes are also reachable by the standard calculus --- $\nu\in S_{45}$,
resp.\ $\nu\in S_{105}$, by explicit four-generator $\pm1$ certificates (re-derivable via the ancillary
lattice test), hence algebraic by Lemma~4.1(iii) of \cite{Aoki00} as well, Theorem~\ref{thm:Aprime}
supplying the explicit two-pair presentation. The computed per-prime values (orders $9/42/7$ at
sampled split primes; computational observations) indicate analogous local residue-degree
constraints at the closed levels; explicit equations remain open.

\section{Methods, code availability, and provenance}\label{sec:methods}
Exact objects throughout (Gauss/Jacobi sums as algebraic integers; the decisive certificate in
$\Z[x]/\Phi_{33}$ with no floating point; the exact-twist certificate framework follows the surface
companion \cite{P1}); numerics corroborate or measure invariant lower bounds only.
Every proof-critical computational claim of Theorem~\ref{thm:B} carries either a stored
per-orbit witness certificate or an algorithmically independent reimplementation; the runner
re-executes the anchor levels live and re-verifies every stored witness and every supplied
certificate against the definitions.
The census engines, the closure oracles, the lattice membership test, the exact Frobenius certificate,
and a standalone verifier of every closure identity in this paper are provided as ancillary files
(\texttt{anc/}; \texttt{run\_all.sh} reproduces the smoke tier --- identities, census anchors, both
lattice verdicts, the Theorem~\ref{thm:C} certificate, the engine calibrations, the
re-verification of all stored per-orbit witnesses including the survivor negatives, the
independent census at anchor levels, and the level-summary cross-check on a fresh sub-range ---
in under a minute on the reference machine whose measured runtimes the ancillary README reports;
\texttt{run\_full\_census.sh} is the archival tier, the
complete from-scratch independent regeneration of all $89$ levels against the pinned summary
hashes; the README records the reference configuration and the measured runtimes quoted in
\S\ref{sec:B}), together with the deep-closure ledger and the stored degree-$210$ certificate as data
artifacts; an algorithmically independent census reimplementation is bundled --- its complete $89$-level
run ships as a checksummed receipt, and the runner re-executes it live at the anchor levels ---
and the original implementation used during verification is methodologically
independent of both. This manuscript was prepared with AI assistance; all
computational claims are backed by the ancillary code and receipts, and the author takes full
responsibility for all proofs, code, and bibliographic claims.

\appendix
\section{The cited Aoki statements}\label{app:aoki}

For self-containedness we reproduce the exact statements used, in the numbering of the primary
texts.

\emph{From \cite{Aoki87}} (notation as there: $\mathfrak B^n_m$ the Hodge characters of $X^n_m$,
$\mathfrak D^n_m$ the classes that are coordinate permutations of $(a_0,-a_0,\dots,a_r,-a_r)$,
$r=n/2$; $*$ juxtaposition; \textsc{claim}$(\alpha)$: $\V(\alpha)$ is generated by classes of
algebraic cycles on $X^n_m$):

\begin{itemize}
\item[--] \textbf{Theorem 1-1} (Shioda). \emph{The linear space $L$ represents
$\delta\in\mathfrak D^n_m$; more precisely
$\omega_\delta(L)\cdot\overline{\omega_\delta(L)}=(-1)^r m^{n+1}$.}
\item[--] \textbf{Theorem 1-4} (Shioda, Ran). \emph{Let $r$ and $s$ be non-negative even integers
such that $n=r+s+2$, and let $\alpha\in\mathfrak B^r_m$, $\beta\in\mathfrak B^s_m$. Then:
(i) if \textup{\textsc{claim}}$(\alpha)$ and \textup{\textsc{claim}}$(\beta)$ are true, then
\textup{\textsc{claim}}$(\alpha*\beta)$ is also true; (ii) if there exists $\delta\in\mathfrak D^s_m$ such
that \textup{\textsc{claim}}$(\alpha*\delta)$ is true, then \textup{\textsc{claim}}$(\alpha)$ is also true.}
\item[--] \textbf{Theorem 2-1.} \emph{For an odd prime $p\mid m$, $d=m/p$, $\gcd(a,d)=1$, the
variety $Y$ defined by (2.1) there is a subvariety of $X^{p-1}_m$ of codimension $r=(p-1)/2$
representing $\alpha=\sigma_{p,a}$; more precisely
$\omega_\alpha(Y)\cdot\overline{\omega_\alpha(Y)}=(-1)^r p^{p-2}m^p$.}
\item[--] \textbf{\S1 definition.} \emph{For a prime $p\mid m$, $d=m/p$, and each $i$ with
$d/(i,d)>2$: $\sigma_{p,i}=(i,\,i+d,\dots,i+(p-1)d,\;m-pi)$ for $p\ge3$ --- the ($p$-)standard
elements.}
\end{itemize}

\emph{From \cite{Aoki00}} (verified verbatim against the open-access copy in the Rikkyo
University repository, doi:10.14992/00009819; \S\ref{sec:even} records two errata of the printed
text, both likewise verified against the primary --- the $\xi_{28}$ list entry on p.~185, and the
attribution line after Theorem~1.2 of the companion CMUSP \textbf{51} (2002) paper, which assigns
its part (ii) (= Theorem~0.1(i) below) to the Catalan-curves reference):

\begin{itemize}
\item[--] \textbf{Theorem 0.1.} \emph{Suppose the prime factorization of $m$ is one of the
following forms: (i) $m=2^a3^b5^c7^d$, where $a,b,c,d$ are non-negative integers such that either
$c=0$ or $d=0$; (ii) $m=p^e$ or $2p^e$, $p$ an odd prime. Then the Hodge conjecture is true for
all abelian varieties of Fermat type of degree $m$.} Its proof (\S4 there) opens: ``In view of
Corollary~3.2, it suffices to show that the Hodge conjecture for $X^n_m$ is true for all $n$'' ---
the bridge to the Fermat varieties themselves used in \S\ref{sec:even}; the
verification then reduces, via Lemma~4.1 and Theorem~2.1 there, to the gap generators $\xi_m$ at
$m=12,15,20,21,28$ (the $m=28$ entry being the misprint recorded in \S\ref{sec:even}).
\item[--] \textbf{Lemma 4.1} (verbatim). \emph{Let $\alpha,\beta\in B_m$. Then:
(i) if both $\V(\alpha)$ and $\V(\beta)$ are algebraic, then so is $\V(\alpha*\beta)$;
(ii) if $\V(\alpha*\delta)$ is algebraic for some $\delta\in D_m$, then so is $\V(\alpha)$;
(iii) if $\alpha\in S_m$, then $\V(\alpha)$ is algebraic.} No hypothesis on $m$ is imposed; the
groups $S_m$ and the ``gap group'' are Aoki's own objects (\S2 there), and
Proposition~\ref{prop:ident} identifies the implemented lattice with $S_m$
generator-by-generator (his \S2 also proves, via Theorem~2.1, that $2\alpha\in S_m$ for every
$\alpha\in B_m$ --- the doubling our $2u$-certificates instantiate).
\end{itemize}

\begin{thebibliography}{99}
\bibitem{Aoki87} N.~Aoki, \emph{Some new algebraic cycles on Fermat varieties}, J.~Math.\ Soc.\ Japan
\textbf{39} (1987), 385--396.
\bibitem{Aoki00} N.~Aoki, \emph{Some remarks on the Hodge conjecture for abelian varieties of Fermat
type}, Comment.\ Math.\ Univ.\ Sancti Pauli \textbf{49} (2000), 177--194; doi:10.14992/00009819
(open access, Rikkyo University repository).
\bibitem{daSilva} G.~da Silva Jr., \emph{Notes on the Hodge conjecture for Fermat varieties},
Experimental Results \textbf{2} (2021), e22, doi:10.1017/exp.2021.14; arXiv:2101.04739.
\bibitem{P1} R.~Jumagulov, \emph{Galois-invariant N\'eron--Severi ranks of Fermat surfaces over
number fields: a Galois module, closed forms, a threshold, and exact tables}, arXiv:2607.17387.
\bibitem{BNeven} R.~Jumagulov, \emph{Residual gap classes in the even-degree Fermat-fourfold
census through degree 250: exchange walls and two closures}, companion paper, 2026 (submitted
alongside; the ancillary package is shared).
\bibitem{Markman1} E.~Markman, \emph{Cycles on abelian $2n$-folds of Weil type from secant sheaves
on abelian $n$-folds}, arXiv:2502.03415.
\bibitem{Markman2} E.~Markman, \emph{Secant sheaves on abelian $n$-folds with real multiplication
and Weil classes on abelian $2n$-folds with complex multiplication}, arXiv:2509.23079.
\bibitem{Schoen} C.~Schoen, \emph{Hodge classes on self-products of a variety with an automorphism},
Compositio Math.\ \textbf{65} (1988), 3--32; Addendum, \textbf{114} (1998), 329--336.
\bibitem{Shioda79} T.~Shioda, \emph{The Hodge conjecture for Fermat varieties}, Math.\ Ann.\ \textbf{245}
(1979), 175--184.
\bibitem{SK} T.~Shioda, T.~Katsura, \emph{On Fermat varieties}, T\^ohoku Math.\ J.\ \textbf{31} (1979),
97--115.
\bibitem{Weil49} A.~Weil, \emph{Numbers of solutions of equations in finite fields}, Bull.\ Amer.\
Math.\ Soc.\ \textbf{55} (1949), 497--508.
\bibitem{Weil52} A.~Weil, \emph{Jacobi sums as ``Gr\"ossencharaktere''}, Trans.\ Amer.\ Math.\ Soc.\
\textbf{73} (1952), 487--495.
\end{thebibliography}
\end{document}
